\documentclass[../main-en.tex]{subfiles}
\begin{document}

\section{Allocation View}
\label{sec:allocation-view}

Allocation structures map software elements to things that are not software:
files, hardware, teams~\cite{bass_software_2021}. The hub is a single
process on one machine, so this view is short, but it makes the cost of a
change measurable (\cref{tab:allocation-changes}). A work
assignment view is left out; for a project of this size it would bring
nothing.

\subsection{Implementation}

Each module of the module view is one C\# project, built into one assembly;
the \texttt{manifest} dependency in \cref{fig:allocation-modules} says that
the assembly contains the code of the module. In .NET~8 every project,
including the entry project, is built into a \texttt{.dll}, with a small
launcher next to the entry assembly.

\landscapefigure{uml/allocation-view/Modules.png}{Implementation view: each module and the assembly built from it.}{fig:allocation-modules}

\Cref{fig:allocation-references} shows the project references between the
assemblies; a \texttt{reference} dependency is a \texttt{<ProjectReference>}
in the \texttt{.csproj} file. The graph is exactly the uses view of
\cref{sec:module-view}, drawn in the same layout. And a project can only
use the types of the projects it references, so a forbidden dependency does
not compile: the rules of the uses view are checked by the compiler on
every build.

\landscapefigure{uml/allocation-view/References.png}{Implementation view: project references between the assemblies,
  same graph as the uses view.}{fig:allocation-references}

We first drew one implementation diagram with packages, assemblies,
manifests and references all together; the use arrows and the reference
arrows crossed and the result was unreadable, so we split it into the two
diagrams above and gave the references diagram the layout of the uses view.

\subsection{Deployment}

\Cref{fig:allocation-deployment} shows the deployment: a laptop with the
.NET~8 runtime installed as an execution environment, the six
assemblies and the data files nested in it (nesting is the UML notation for
``deployed on''). The console dashboard is not drawn: it is the terminal
where the hub is started.

\landscapefigure{uml/allocation-view/Deployment.png}{Deployment view: the laptop, the .NET~8 runtime, the assemblies
  and the files.}{fig:allocation-deployment}

The first program argument is the path of the configuration file
(\texttt{config.json} if there is none), and relative paths are resolved
against the output folder rather than the current directory, so the hub
finds its files whatever directory it is started from.

\subsection{Impact of likely changes}

\Cref{tab:allocation-changes} reads the cost of a change off the project
references: when a project changes, the projects that reference it are
recompiled too.

\begin{table}[H]
  \centering
  \small
  \caption{Impact of likely changes across the allocation view.}
  \label{tab:allocation-changes}
  \begin{tabularx}{\textwidth}{| L | L | L | L |}
    \hline
    \textbf{Change} & \textbf{Edited} & \textbf{Recompiled} &
    \textbf{Redeployed} \\
    \hline
    Chain order, conversion, threshold, window size & \texttt{config.json} &
    none & nothing, restart the hub \\
    \hline
    Add a transformation & one class in Processing, one entry in the
    transformation factory, \texttt{config.json} & Processing,
    Configuration, Application & the output folder \\
    \hline
    Add a source & one class in Ingestion, one entry in the
    \texttt{SourceFactory}, \texttt{config.json} & Ingestion, Configuration,
    Application & the output folder \\
    \hline
    Add a sink & one class in Output, one entry in \texttt{SinkFactory},
    \texttt{config.json} & Output, Configuration, Application & the output
    folder \\
    \hline
    Change a shared data type in Core & Core and every module that uses the
    changed member & all six & the output folder \\
    \hline
  \end{tabularx}
\end{table}

The most frequent changes need no compilation at all, because the chain, the
conversions and the thresholds live in \texttt{config.json}, a file next to
the binaries~\cite{wiklund_ch8_nodate}. Adding a source, transformation or
sink never touches the other two kinds, and Core is the expensive one: every
project references it, which is a second reason to keep it to plain data
types.

\end{document}
